\documentclass[12pt,a4paper]{article}

\usepackage[italian]{babel}
\usepackage{geometry}
\geometry{margin=2.5cm}
\usepackage{array}
\usepackage{tabularx}
\usepackage{xcolor}

\definecolor{negativeRed}{RGB}{180, 40, 40}
\definecolor{orangino}{RGB}{205, 127, 50}

\newcommand{\cognitivetable}[4]{%
	\begin{center}
		\renewcommand{\arraystretch}{1.5}
		\begin{tabularx}{\linewidth}{|X|X|}
			\hline 
			\textbf{Domanda} & \textbf{Osservazioni}
			\\ \hline
			\textbf{1. L'utente sa cosa deve ottenere con l'azione?} (sa cosa deve fare?) & #1 \\ \hline
			\textbf{2. L'utente può individuare lo strumento di interazione?} & #2 \\ \hline
			\textbf{3. L'utente può anche capire che quello è lo strumento giusto per fare ciò che vuol fare?} & #3 \\ \hline
			\textbf{4. Dopo che l'azione è eseguita, l'utente capisce la risposta che ottiene?} & #4 \\ \hline
		\end{tabularx}
	\end{center}
}

\begin{document}
	\section*{Documento di valutazione}
	
	\noindent
	\textbf{Codice di valutazione:} 2 \\
	\textbf{Valutatore:} Singh Sahiljit  \\
	\textbf{Data di valutazione:} 03/08/2026 \\
	\textbf{Sistema da valutare:} Sito web Home Banking \\
	\textbf{Scopo:} Analizzare le difficoltà dell'utente nell' usare il sistema\\
	\textbf{Compito:} Fare un bonifico\\
	
	\subsection*{Valutazione Azione}
	
	Azione 1: Passare alla sezione per fare un bonifico.
	\cognitivetable{%
		Sì è chiaro.
	}{%
		Sì il bottone "Bonifio" è chiaramente visibile %
	}{%
		Sì in quanto non ci sono ambiguità%
	}{%
		Sì. Il sistema mostra la pagina per eseguire il bonifico%
	}
	
	\noindent
	Azione 2: Inserire i dati del beneficiario.
	\cognitivetable{%
		Sì è chiaro.
	}{%
		Sì il forum per l'input è evidente%
	}{%
		Sì in quanto non ci sono ambiguità%
	}{%
		\textcolor{orangino}{Sì e no. Il testo che conferma la correttezza dei dati inseriti è illeggibile oppure non si nota}%
	}

	\clearpage
	\noindent
	Azione 3: Inserire l’importo e causale.
	\cognitivetable{%
		Sì è chiaro.
	}{%
		Sì il forum per l'input è evidente%
	}{%
		Sì in quanto non ci sono ambiguità%
	}{%
		Sì appare un elemento visivo che fa capire all'utente che ha svolto tutte le operazioni di inserimento correttamente%
	}
	
	\noindent
	Azione 4: Confermare o annullare il bonifico.
	\cognitivetable{%
		Sì è chiaro.
	}{%
		Sì il bottone per confermare diventa abilitato. Per quanto riguarda l'annullamento anche questo bottone è chiaramente visibile vicino al bottone di conferma con una certa spaziatura fra i due.%
	}{%
		Sì in quanto non ci sono ambiguità%
	}{%
		Sì appare una schermata che mostra che il bonifico è andato a buon fine%
	}
\end{document}